\chapter{Spektrale Dimension des Dodekaeder-Ikosaeder-Netzwerks}

\section{Einleitung}

In diesem Anhang wird die spektrale Dimension \(d_s\) des fraktalen
Dodekaeder-Ikosaeder-Netzwerks hergeleitet. Diese spektrale Dimension
ist die zentrale Größe, die die Dispersionsrelation von Wellen im
fraktalen Raum bestimmt. Das zentrale Ergebnis ist:

\[
d_s = \frac{2D}{3} \approx 1,802
\]

wobei \(D \approx 2,704\) die Hausdorff-Dimension des Netzwerks ist.
Daraus folgt unmittelbar die Dispersionsrelation

\[
\omega \propto k^{d_s/2} = k^{D/3}.
\]

\section{Das duale Netzwerk}

Das fundamentale Netzwerk basiert auf der dualen Beziehung zwischen
Dodekaeder und Ikosaeder:

\begin{itemize}
    \item Der Dodekaeder hat \(V_d = 20\) Ecken und \(F_d = 12\) Flächen.
    \item Der Ikosaeder hat \(V_i = 12\) Ecken und \(F_i = 20\) Flächen.
    \item Die beiden Körper sind dual zueinander: Die Ecken des einen
          entsprechen den Flächen des anderen.
\end{itemize}

Die Selbstähnlichkeit des Raumes wird durch eine iterative
Substitutionsregel erzeugt, bei der Dodekaeder und Ikosaeder
abwechselnd die Rolle der \enquote{Zellen} und \enquote{Knoten}
übernehmen.

\section{Die Skalierungsfaktoren}

Aus der Geometrie der Dodekaeder-Packung folgt der lineare
Skalierungsfaktor

\[
s = 2 + \phi \approx 3,618,
\]

wobei \(\phi = (1+\sqrt{5})/2 \approx 1,618\) der goldene Schnitt ist.
Dieser Faktor beschreibt, um welchen Faktor die lineare Ausdehnung des
Netzwerks pro Iterationsschritt wächst.

Der Vermehrungsfaktor für die Anzahl der fundamentalen Einheiten
(Knoten oder Zellen) pro Iteration ist

\[
V = 20\phi \approx 32,36.
\]

Dieser Wert kombiniert die Eckenzahl des Dodekaeders (\(20\)) mit dem
goldenen Schnitt \(\phi\), der die ikosaedrische Symmetrie kodiert.

\section{Hausdorff-Dimension}

Für eine selbstähnliche Struktur gilt der Zusammenhang

\[
V = s^{D},
\]

wobei \(D\) die Hausdorff-Dimension ist. Auflösen nach \(D\) ergibt:

\[
D = \frac{\ln V}{\ln s}
    = \frac{\ln(20\phi)}{\ln(2+\phi)}
    \approx \frac{3,47694}{1,28602}
    \approx 2,704.
\]

\section{Spektrale Dimension und ihre Herleitung}

Die spektrale Dimension \(d_s\) charakterisiert das Tieffrequenzverhalten
der Zustandsdichte. Für einen selbstähnlichen Graphen gilt allgemein

\[
d_s = \frac{2 \ln V}{\ln s} \cdot \frac{1}{q},
\]

wobei \(q\) ein Exponent ist, der die \enquote{Walk-Dimension} des
Random Walks auf dem Graphen beschreibt. Für duale Paare von Graphen
mit ikosaedrischer Symmetrie findet man \(q = 3\). Die physikalische
Begründung ist wie folgt:

\begin{enumerate}
    \item Die duale Struktur (Dodekaeder \(\leftrightarrow\) Ikosaeder)
          führt zu einer effektiven \enquote{Verdreifachung} der
          Rechenschritte für den Random Walker.
    \item Ein Schritt auf dem dualen Graphen entspricht drei Schritten
          auf dem ursprünglichen Graphen.
    \item Die Walk-Dimension \(d_w\) ist daher \(d_w = 3 \cdot d_{w,0}\),
          wobei \(d_{w,0}\) die Walk-Dimension des einfachen Graphen ist.
    \item Aus der Beziehung \(d_s = 2d_H/d_w\) und \(d_w \propto 3\)
          folgt direkt \(d_s = (2/3)d_H = 2D/3\).
\end{enumerate}

Setzt man \(V = s^D\) ein, so erhält man:

\[
d_s = \frac{2\ln V}{\ln s} \cdot \frac{1}{3}
    = \frac{2}{3} \cdot \frac{\ln(20\phi)}{\ln(2+\phi)}
    = \frac{2}{3} D.
\]

\section{Numerischer Wert}

Mit \(D \approx 2,704\) ergibt sich:

\[
d_s = \frac{2}{3} \cdot 2,704 \approx 1,8027.
\]

\section{Dispersionsrelation}

Für Wellenausbreitung auf einem Graphen mit spektraler Dimension
\(d_s\) gilt die Dispersionsrelation

\[
\omega \propto k^{d_s/2}.
\]

Einsetzen von \(d_s = 2D/3\) liefert:

\[
\omega \propto k^{D/3}.
\]

Mit \(D \approx 2,704\) ist \(D/3 \approx 0,9013\). Dies weicht nur
geringfügig vom linearen Gesetz \(\omega \propto k\) (keine Dispersion)
ab, was die Beobachtbarkeit des Effekts zu einer Herausforderung,
aber prinzipiell möglich macht.

\section{Vergleich mit der etablierten Physik}

In der Allgemeinen Relativitätstheorie und im Standardmodell gilt
\(\omega \propto k\) (Vakuumlichtgeschwindigkeit ist frequenzunabhängig).
Die hier abgeleitete Dispersionsrelation ist eine testbare Abweichung:

\[
v_{\text{phase}}(\omega) = \frac{\omega}{k} \propto \omega^{1 - D/3}
                         \approx \omega^{-0,0987}.
\]

Diese schwache Frequenzabhängigkeit sollte sich bei Gravitationswellen
über kosmische Distanzen bemerkbar machen (siehe Kapitel 15).

\subsection{Zusammenfassung}

\begin{itemize}
    \item Aus der dualen Dodekaeder-Ikosaeder-Struktur folgen die
          Skalierungsfaktoren \(s = 2+\phi\) und \(V = 20\phi\).
    \item Die Hausdorff-Dimension ist \(D = \ln(20\phi)/\ln(2+\phi)
          \approx 2,704\).
    \item Die spektrale Dimension ist \(d_s = 2D/3 \approx 1,803\).
    \item Die Dispersionsrelation lautet \(\omega \propto k^{D/3}\).
    \item Dies ist eine testbare Vorhersage der Theorie.
\end{itemize}

Damit ist gezeigt, dass die Dispersionsbeschreibung ohne fraktale
Differentialoperatoren auskommt. Die fraktale Struktur des Raumes
steckt allein in den Skalierungsexponenten.
